\chapter{可实施特征与实施例}
\label{ch:patent}

前十章的判据在装置上收成三组可分开实施的步骤。三组共享符号、共享执行器，分别对应信息获取、路径解释和安全执行。

\section{三组步骤的分界}

\begin{table}[H]
\centering
\caption{三组步骤各自必不可少的特征}
\begin{tabularx}{\textwidth}{lXX}
\toprule
组别 & 必不可少的特征 & 不依赖的对象 \\
\midrule
信息组 & 相变变号脉冲后接对流脉冲；协方差未下降则不改停止时刻 & 具体食谱、网络结构 \\
解释组 & 冷点路径与表面路径分别积分；符号不稳定则只保留区间 & 从数据中搜索父节点 \\
安全组 & 投影只用硬测量；断电回路独立；低置信度改为保温轨迹 & 屏障理论的一般形式 \\
\bottomrule
\end{tabularx}
\end{table}

跨次参数写回属于信息组的后续步骤：探针有效且精度高于先验才写入本地类别分布。单独的“在线修正烹饪曲线”不包含无效探针否决，也不包含精度门槛。

\section{信息组}

非均质食品热加工状态的辨识与控制，按下列顺序实施。

\begin{enumerate}
\item 状态包含温度、水分、反应进度与结构；决策量取冷点热力进度、表面温度、表面含水与内部热阻。
\item 由决策量协方差的行列式或 Fisher 信息判断信息赤字。
\item 赤字高于阈值时，蒸汽执行器产生一段使界面相变率变号的第一脉冲；排汽与气流执行器随后产生一段使外部对流热阻下降、界面转为蒸发的第二脉冲。
\item 用两段期间表面温度、腔体湿度和食品质量之差，更新内部热阻与热容。
\item 协方差下降未达到预定比例，则该次脉冲无效，禁止据此改变停止时刻。
\item 下降达到预定比例，后验才进入后续执行量计算。
\end{enumerate}

可并入的限制：脉冲时长上限、表面褐变增量预算、质量通道缺失时否决、厚度超过信息上限时输出区间。

\section{解释组}

热加工预测按下列顺序给出路径，而不是给出单一剩余时间。

\begin{enumerate}
\item 结构方程按腔体水分、界面相变、内部导热、反应进度、结构到热阻的顺序固定父节点。
\item 对同一执行器，沿潜热至冷点进度、表面含水至表面品质至少两条路径分别积分。
\item 替换该执行器轨迹并保持其余方程，得到各路径对品质泛函的偏效应和后验区间。
\item 偏效应符号在后验抽样中一致，才输出该路径的语句；不一致则只输出预测区间。
\end{enumerate}

反事实使用已观测窗口反推的外生扰动，界面同时给出两条路径的数值。

\section{安全组}

执行按下列顺序约束品质模块。

\begin{enumerate}
\item 品质模块给出候选执行量。
\item 投影只读取漏电相关量、水位或空烧温升、加热器温度、腔压与门状态，求与候选量距离最小且满足屏障不等式的执行量。
\item 漏电、干烧或超温另有不经过品质模块处理器的断电回路。
\item 协方差或残差越界时，候选量改为存储的低功率保温轨迹，再进入同一次投影。
\end{enumerate}

探针期间同样投影。投影超时则维持上一周期已经投影的执行量。

\section{在同一腔体上的实施例}

蒸汽、排汽、风机和加热器既产生加工边界，也产生信息组的两段脉冲。质量与热像是两段脉冲的采样。解释组的路径差显示为冷点项与表面项两个数。安全组的前半段是投影，后半段是剩余电流、水位和管温的硬件切断。

腔体把蒸与烤的边界放在一起，只说明执行器已经具备。可分开实施的是脉冲顺序、路径分解和投影所用的测量集合。

\begin{figure}[H]
\centering
\begin{tikzpicture}[
  box/.style={draw, rounded corners, align=center, font=\small, inner sep=3pt, minimum height=8mm},
  arr/.style={-{Stealth}}
]
\node[box] (s) {蒸汽 / 排汽 / 风 / 加热};
\node[box, right=8mm of s] (f) {界面相变\\食品内部};
\node[box, right=8mm of f] (y) {质量、热像、湿度};
\node[box, below=8mm of y] (p) {后验与路径差};
\node[box, below=8mm of s] (h) {漏电、水位、管温、门};
\node[box, right=8mm of h] (u) {投影后的边界};
\draw[arr] (s)--(f);
\draw[arr] (f)--(y);
\draw[arr] (y)--(p);
\draw[arr] (p.west) -- ++(-6mm,0) |- (u.east);
\draw[arr] (h)--(u);
\draw[arr] (u.north) -- ++(0,4mm) -| (s.south);
\end{tikzpicture}
\caption{实施例中的两条测量链。上方链进入后验，下方链只进入投影。}
\end{figure}

\section{结果尚未产生时的文档边界}

机制、判据和试验协议在前十章已经闭合。真值模型尚未运行的网格，不出现曲线、百分比和剩余时间的示例数字。信息上限、路径符号和越界次数以第 \ref{ch:sim} 章的指标定义为准，由 \texttt{eval/} 从真值日志计算。

\section{实施顺序}

\begin{enumerate}
\item 冻结两节点方程、观测方程、P4 时序和指标脚本。
\item 用一维真值模型完成 T 与 E 不同的四策略对照，先不接入神经算子。
\item 加入 Fisher 打分，确认有效探针的 $G_{\mathrm{real}}>0$。
\item 训练算子，在蒸汽干预上与直接回归剩余时间的模型比较。
\item 接入投影和降级，统计越界次数。
\item 三组步骤使用同一套符号写入实施例。
\end{enumerate}

某一步的判据若不成立，收缩对应步骤的适用条件。适用条件来自厚度、噪声和通道缺失，不来自放宽指标。
